% /Users/pikachu/books/payments-book/src/parts/part3/ch17-unionpay-regional.tex
\chapter{UnionPay и региональные карточные схемы}
\label{ch:unionpay}

\begin{kaobox}[title=Что вы узнаете из этой главы]
  Почему UnionPay выбрала режим single message и~какие ограничения
  из~этого вырастают для~банка, мерчанта и~процессинга; как работает
  кобейджинг \scheme{Мир}~\texttimes{}~UnionPay и~почему он~не~отменяет
  географические и~санкционные ограничения; чем национальные схемы
  СНГ похожи на~\scheme{Visa} и~\scheme{Mastercard}, а~в~чём живут
  по~собственной логике~--- в~управлении, процессинге и~трансграничном
  приёме.
\end{kaobox}

Представим карту, которая внутри одной страны работает как~обычный
национальный инструмент, а~за~её пределами должна внезапно
\enquote{переключиться} на~другую сеть. Пока мы~смотрим только
на~логотипы, задача кажется маркетинговой. Но~как только мы~смотрим
на~маршрут сообщения, клиринг и~географию приёма, выясняется, что речь
идёт о~другой архитектуре. Эта глава отвечает на~вопрос, где
региональные схемы повторяют инварианты карточного мира, а~где
сознательно от~них отходят.


\section{UnionPay: почему здесь важен режим single message}
\label{sec:up-architecture}

По~числу карт в~обращении UnionPay~--- крупнейшая карточная схема
в~мире, но~большая часть её~трафика остаётся внутрикитайской
и~идёт через инфраструктуру CUPS (China UnionPay
System)~\todocite{UnionPay Annual Report или официальная спецификация
  CUPS}. Для~такой сети центральная инженерная задача формулируется
иначе, чем~для~\scheme{Visa} или~\scheme{Mastercard}: как упростить
путь между одобрением операции и~её расчётом, если основная масса
транзакций живёт внутри одного большого внутреннего контура.

Ответом становится режим single message. В~классической модели
dual message авторизация и~клиринг разведены во~времени: сначала
эмитент резервирует сумму, потом сеть и~эквайер присылают отдельную
клиринговую запись. В~CUPS авторизационное сообщение несёт
в~себе и~признак одобрения, и~данные для~расчёта. Для~мерчанта
и~банка это~снижает число точек рассогласования: нет отдельного
клирингового файла, меньше причин искать несоответствие между
\enquote{было одобрено} и~\enquote{попало в~расчёт}. Но~упрощение
имеет цену. Гибкость, привычная гостиницам, авиакомпаниям
и~другим сценариям с~pre-auth и~completion, становится заметно
ниже~\todocite{UnionPay technical specifications CUPS}.

QuickPass дополняет эту модель на~уровне пользовательского сценария.
Это бесконтактный сервис UnionPay, работающий поверх ISO/IEC~14443
и~решающий ту~же задачу, что payWave или~PayPass: сократить число
действий покупателя на~кассе. Снаружи покупатель видит знакомый
жест tap-and-pay, но~внутри него действует схема с~другим балансом
между простотой обработки и~последующей гибкостью
операции~\todocite{QuickPass technical specification UnionPay}.

Локальный вывод здесь такой: UnionPay остаётся карточной схемой,
но~выбирает другой компромисс между простотой и~управляемостью.
Если \scheme{Visa} и~\scheme{Mastercard} платят сложностью
последующего reconciliation ради большей гибкости жизненного цикла,
то UnionPay платит этой гибкостью ради более прямого пути
к~расчёту.


\section{Кобейджинг \scheme{Мир} \texttimes{} UnionPay}
\label{sec:up-mir}

Кобейджинг нужен там, где одной национальной инфраструктуры уже
недостаточно, а~возвращение к~глобальной схеме невозможно или
нежелательно. Для~карты \enquote{\scheme{Мир}~--- UnionPay} это
означает простую для~клиента, но~непростую для~процессинга идею:
внутри одного пластика живут два набора правил, и~маршрут операции
зависит от~того, где именно карта предъявлена.

За~пределами России операция по~такой карте идёт по~каналу
UnionPay International, а~не~через НСПК. Терминал, принимающий
UnionPay, считывает совмещённый продукт как карту, которую можно
маршрутизировать через инфраструктуру UnionPay; эмитент получает
авторизационный запрос из~этого контура и~должен уметь на~него
ответить. На~словах это~похоже на~универсальный мост между
внутренним и~внешним миром, но~на~практике мост оказывается уже:
не~все банки подключены к~необходимым международным каналам,
не~все терминалы поддерживают UnionPay, а~часть эквайеров вообще
не~готова работать с~российским риском и~санкционным
контуром~\todocite{актуальный охват Мир-UnionPay по~странам 2024}.

Поэтому кобейджинг не~отменяет ограничений, а~лишь перераспределяет
их. Он~снижает зависимость от~одной-единственной внешней сети,
но~усиливает зависимость от~географии приёма и~от~комплаенс-политики
конкретных банков. Для~архитектора платёжной системы из~этого
следует важный практический вывод: кобейджинг нельзя рассматривать
как~магическую совместимость. Это всегда отдельный маршрут
с~собственными правилами доступности, онбординга и~операционного
риска.


\section{Региональные схемы СНГ}
\label{sec:cis-payment-systems}

После распада СССР многие страны пришли к~одному и~тому~же вопросу:
что произойдёт с~внутренними расчётами, если они целиком завязаны
на~внешнюю инфраструктуру? Ответом стали национальные карточные
схемы, которые должны были удержать внутренний трафик, локальные
комиссии и~операционное управление внутри страны.

Армянская ArCa строилась как национальная схема для~внутреннего
приёма, но~с~возможностью кобейджинга для~международного
использования~\todocite{ArCa national payment system Armenia}.
Белкарт в~Беларуси решает схожую задачу: внутренний приём
и~собственное правило игры внутри страны, при~этом отдельные продукты
связываются с~\scheme{Мир} для~работы в~российском
контуре~\todocite{Белкарт национальная платёжная система Беларуси}.
Узбекистан интересен тем, что там долгое время сосуществовали
две инфраструктуры~--- UzCard и~HUMO,~--- и~сама задача интероперабельности
между ними стала отдельной операционной темой для~банков
и~мерчантов~\todocite{HUMO UzCard payment systems Uzbekistan}.
ЭЛКАРТ в~Кыргызстане подчёркивает ещё одну грань таких систем:
национальная схема может быть построена прежде всего вокруг
внутренних выплат и~государственных сценариев, а~не~вокруг
международного travel-use case~\todocite{ЭЛКАРТ Кыргызстан}.

\todofig{Карта СНГ с~национальными платёжными системами.
  Карта постсоветского пространства. Над каждой страной~---
  подпись с~названием схемы: Россия~--- «Мир (НСПК)»;
  Беларусь~--- «Белкарт + Мир»; Армения~--- «ArCa»;
  Узбекистан~--- «HUMO + UzCard»; Кыргызстан~--- «ЭЛКАРТ»;
  Казахстан~--- «локальные проекты + Visa/MC»; Украина~--- «ПРОСТІР».
  Цветом показать степень опоры на~собственную внутреннюю
  инфраструктуру: тёмный тон~--- основной внутренний rail,
  светлый тон~--- сильная зависимость от~кобейджинга.
  Стиль: минималистичная политическая карта без~декоративных
  элементов.}
% \label{fig:cis-payment-map}

Здесь важно не~список названий, а~общий инвариант. Почти все такие
системы технически опираются на~EMV и~\iso{8583}, то~есть говорят
на~языке, знакомом международным картам. Но~жизнь схемы определяется
не~только форматом сообщения. Её определяют управление системой,
правила допуска участников, география обязательного приёма,
кобейджинговые договорённости и~пределы трансграничного использования.
Именно поэтому региональная схема может быть очень похожа на~глобальную
по~формату пакета и~очень непохожа по~операционной логике.

Когда эта граница установлена, можно двигаться дальше. В~региональных
карточных схемах карта всё ещё остаётся центром модели, даже если
правила приёма и~маршрутизации отличаются от~международных сетей.
Следующая глава задаёт более радикальный вопрос: что меняется, когда
центром становится уже не~карта, а~счёт и~доступ к~нему через API.


\section*{Итоги главы}

\begin{itemize}
  \item UnionPay остаётся карточной схемой, но~выбирает режим
        single message, уменьшая число рассогласований между
        одобрением и~расчётом ценой меньшей гибкости жизненного цикла.
  \item QuickPass и~кобейджинг показывают, что внешне знакомый
        карточный сценарий может опираться на~другую внутреннюю
        процессинговую логику.
  \item Национальные схемы СНГ обычно используют EMV и~\iso{8583},
        но~живут по~собственным правилам управления, приёма и~географии.
  \item Кобейджинг и~региональные схемы не~отменяют трансграничные
        ограничения, а~лишь по-другому распределяют инфраструктурный
        и~комплаенс-риск.
\end{itemize}
